% ============================================================
% CS 5365 Final Project -- 2-Page Summary Report (excl. references)
% Re-implementing Attention Is All You Need
% Authors: Arman, Samin Islam
% ============================================================
\documentclass[10pt,twocolumn]{article}

\usepackage[top=0.55in,bottom=0.55in,left=0.55in,right=0.55in,columnsep=0.25in]{geometry}
\usepackage{graphicx}
\usepackage{booktabs}
\usepackage{caption}
\usepackage{hyperref}
\usepackage{amsmath}
\usepackage{microtype}
\usepackage{enumitem}
\usepackage{titlesec}
\usepackage{xcolor}
\usepackage{parskip}

\definecolor{navy}{HTML}{1E3A5F}
\definecolor{dgray}{HTML}{475569}

\titleformat{\section}
  {\normalfont\normalsize\bfseries\color{navy}}
  {\thesection.}{0.4em}{}[\vspace{-4pt}\textcolor{dgray}{\hrule}\vspace{2pt}]
\titlespacing{\section}{0pt}{6pt}{2pt}

% Custom command for the unnumbered References heading (same style)
\newcommand{\refsheading}{%
  \vspace{4pt}%
  {\normalfont\normalsize\bfseries\color{navy}References}%
  \vspace{-4pt}\par\textcolor{dgray}{\hrule}\vspace{2pt}%
}

\captionsetup{font=small,labelfont=bf,skip=3pt}
\setlist[itemize]{noitemsep,topsep=1pt,parsep=0pt,leftmargin=*}

\hypersetup{colorlinks=true,linkcolor=navy,urlcolor=navy}
\setlength{\parindent}{0pt}
\setlength{\parskip}{3pt}

\begin{document}

% -- Title block (spans both columns) ------------------------
\twocolumn[{%
\begin{center}
  {\large\bfseries\color{navy}
    Re-implementing ``Attention Is All You Need'':\\[2pt]
    A PyTorch Study on Multi30k DE$\rightarrow$EN Translation}\\[4pt]
  {\small
    \textbf{Arman} \quad$\cdot$\quad \textbf{Samin Islam}\\[2pt]
    CS~5365: Deep Learning \quad$\cdot$\quad
    The University of Texas at El~Paso}\\[5pt]
  \textcolor{navy}{\rule{\linewidth}{1.2pt}}
\end{center}
\vspace{2pt}
}]

% ============================================================
\section{Introduction}

\textbf{Problem and motivation.}
Sequence-to-sequence tasks such as machine translation had long depended on
recurrent neural networks (RNNs) or convolutional architectures. These designs
process tokens sequentially, limiting parallelism and making long-range
dependency modeling difficult.

\textbf{Paper being reproduced.}
\textit{Attention Is All You Need} -- Vaswani, Shazeer, Parmar, Uszkoreit,
Jones, Gomez, Kaiser, and Polosukhin (NeurIPS 2017) -- proposed the
\textbf{Transformer}: the first sequence-to-sequence model built
\emph{entirely} on self-attention, with no recurrence or convolution.

\textbf{Primary method and contributions.}
\begin{itemize}
  \item \textbf{Multi-head self-attention:} 8 parallel attention heads each
    learn different relational patterns across all positions simultaneously.
  \item \textbf{Positional encoding:} sinusoidal functions inject sequence
    order without recurrence, enabling full parallelism.
  \item \textbf{Encoder-decoder architecture:} 6 stacked layers each in the
    encoder and decoder, connected via cross-attention.
  \item \textbf{State-of-the-art results:} 28.4 BLEU on WMT14 EN-DE and
    41.0 BLEU on EN-FR, surpassing all prior models at lower training cost.
\end{itemize}

Our goal was to re-implement the base Transformer in PyTorch, modernize a
legacy codebase for Apple Silicon (PyTorch~2.x MPS), and verify that the
reported training dynamics are reproducible on the smaller Multi30k
DE$\rightarrow$EN dataset ($\sim$29\,k sentence pairs).

% ============================================================
\section{Chosen Result}

We targeted the \textbf{training convergence behavior} described in Section~6
of the paper, visualized in \textbf{Figure~2} (training loss curves), for the
\emph{base model} configuration in \textbf{Table~3}.

The specific result is the \textbf{Noam learning-rate schedule}
(\textbf{Equation~3} of the paper):
\[
  \mathrm{lr} = d_{\mathrm{model}}^{-0.5}\!\cdot\!
    \min\!\left(\mathrm{step}^{-0.5},\;
    \mathrm{step}\cdot\mathrm{warmup}^{-1.5}\right)
\]
which linearly warms up the learning rate for the first \texttt{warmup} steps,
then decays it as the inverse square root of the step number.

\textbf{Why this result?}
The schedule is central to the paper's training claim: that an
attention-only model with correct optimization \emph{converges stably}
and outperforms RNN-based systems. Reproducing this dynamic -- even on a
smaller dataset -- directly validates the core training methodology of the
Transformer and demonstrates that the architecture's behavior is not
dataset-specific.

% ============================================================
\section{Methodology}

\textbf{Model architecture.}
We implemented the base Transformer exactly as specified: $d_{\mathrm{model}}
{=}512$, 6 encoder and 6 decoder layers, 8 attention heads,
$d_{\mathrm{ff}}{=}2048$, $d_k{=}d_v{=}64$, dropout~$=0.1$, sinusoidal
positional encodings. Source, target, and pre-softmax projection weights are
shared (Section~3.4 of the paper).

\textbf{Dataset.}
WMT'16 Multi30k DE$\rightarrow$EN ($\sim$29\,k train / $\sim$1\,k validation
pairs), tokenized with spaCy~3.7, shared vocabulary. Multi30k was chosen as
a tractable proxy for WMT14 given hardware constraints; we acknowledge this
limits absolute BLEU comparability.

\textbf{Training setup.}
80 epochs; batch size~256; Adam ($\beta_1{=}0.9$, $\beta_2{=}0.98$,
$\varepsilon{=}10^{-9}$); Noam schedule with \texttt{lr\_mul}$\,{=}\,2.0$,
warmup $= 128{,}000$ steps; label smoothing $\varepsilon{=}0.1$.
Hardware: Apple M4 Pro 24\,GB, PyTorch~2.x MPS backend.

\textbf{Evaluation metrics.}
Token-level \emph{cross-entropy loss}, \emph{perplexity}
($e^{\mathcal{L}}$), and \emph{token accuracy} on both splits per epoch.
BLEU was not computed as beam-search decoding was not activated in this run.

\textbf{Modifications from the original.}
The 2019 codebase (PyTorch~1.3) required four changes:
\begin{itemize}
  \item \textbf{torchtext API:} \texttt{Field}, \texttt{Dataset}, and
    \texttt{BucketIterator} were removed in torchtext~0.9. We replaced them
    with a custom \texttt{compat.py} shim.
  \item \textbf{PyTorch 1.3$\rightarrow$2.x:} required for MPS (Apple
    Silicon GPU) support.
  \item \textbf{spaCy 2$\rightarrow$3.7:} arm64 wheels unavailable in v2.
  \item \textbf{Removed unused packages:} \texttt{msgpack-python},
    \texttt{msgpack-numpy}, \texttt{terminado} (never imported in source).
\end{itemize}

% ============================================================
\section{Results \& Analysis}

\begin{figure*}[t]
  \centering
  \includegraphics[width=\linewidth]{training_results.png}
  \caption{Training dynamics over 80 epochs on Multi30k DE$\rightarrow$EN.
    \textbf{(A)} Perplexity, \textbf{(B)} Token accuracy,
    \textbf{(C)} Cross-entropy loss, \textbf{(D)} Noam LR schedule
    (warmup\,=\,4\,k steps shown). Blue\,=\,Train, Red\,=\,Validation.}
  \label{fig:results}
\end{figure*}

Figure~\ref{fig:results} shows the full 80-epoch curves; Table~\ref{tab:compare}
compares our results against the original paper's reported numbers.

\begin{table}[h]
  \centering
  \caption{Re-implementation vs.\ original paper (base model).}
  \label{tab:compare}
  \small
  \setlength{\tabcolsep}{4pt}
  \begin{tabular}{lccc}
    \toprule
    & \textbf{Paper} & \textbf{Ours (train)} & \textbf{Ours (val)} \\
    \midrule
    Dataset      & WMT14     & Multi30k       & Multi30k \\
    Pairs        & 4.5M      & 29k            & --       \\
    BLEU EN-DE   & 27.3      & N/A            & N/A      \\
    Best PPL     & --        & 7.28 (ep.78)   & 16.02 (ep.27) \\
    Token Acc.   & --        & 79.4\% (ep.78) & 48.3\% (ep.52) \\
    \bottomrule
  \end{tabular}
\end{table}

\textbf{What matched the paper.}
The Noam schedule produced the characteristic LR curve from Equation~3:
linear warm-up followed by inverse-square-root decay (Figure~\ref{fig:results}D).
Training perplexity fell monotonically from $\sim$2627 at epoch~0 to 7.28 at
epoch~78, matching the smooth convergence shown in the paper's Figure~2.
Adam with the paper's exact hyperparameters and label smoothing produced
stable, well-behaved loss curves throughout.

\textbf{Discrepancies and challenges.}
Validation perplexity bottomed at 16.02 (epoch~27) then rose, while training
accuracy climbed to 79.4\% -- a clear overfitting signal. This was
expected: Multi30k is $\sim$150$\times$ smaller than WMT14, leaving the
model insufficient data. BLEU could not be directly compared due to the
absence of beam-search decoding. The main engineering challenge was
replacing the deprecated torchtext API; the custom \texttt{compat.py}
shim required careful handling of padding indices, special tokens
(\texttt{<s>}, \texttt{</s>}, \texttt{<blank>}), and dynamic batching to
exactly match the original data pipeline behavior.

\textbf{Broader analysis.}
Despite overfitting on the small corpus, the results confirm the paper's
central claim: multi-head self-attention alone is sufficient to learn
translation structure. The training dynamics are architecture-driven, not
dataset-specific. The full power of the Transformer emerges at scale
(WMT14 / WMT17), where its parallelism makes training feasible in ways
that sequential RNNs cannot match.

% ============================================================
\section{Reflections}

\textbf{Lessons learned.}
\begin{itemize}
  \item \textbf{Implementation deepens understanding.} Coding scaled
    dot-product attention clarified \emph{why} $\sqrt{d_k}$ scaling is
    necessary: without it, large dot products saturate the softmax and
    produce near-zero gradients. The Noam schedule similarly made clear
    why warm-up prevents divergence at random initialization.
  \item \textbf{The schedule is non-negotiable.} Early test runs without
    warm-up diverged immediately, confirming the Noam schedule is a core
    ingredient, not an optional improvement.
  \item \textbf{Infrastructure ages fast.} A three-year-old codebase
    required rewriting its entire data pipeline. Long-term reproducibility
    demands pinned environments and thorough dependency documentation --
    not just published model weights.
  \item \textbf{Scale is not a detail.} The gap between Multi30k and WMT14
    explains nearly all discrepancy with the paper. Dropout and label
    smoothing help but cannot substitute for data volume at this model size.
\end{itemize}

\textbf{Future directions.}
\begin{itemize}
  \item Train on WMT14 EN-DE to reproduce the 27.3 BLEU base-model result
    from Table~3 of the paper.
  \item Enable beam-search decoding and compute sacreBleu for a direct
    comparison against the paper's reported scores.
  \item Visualize individual attention heads to study what linguistic
    relationships (e.g., coreference, subject-verb agreement, syntax) each
    of the 8 heads captures.
  \item Implement the \emph{big} Transformer variant to validate the
    scaling results also reported in Table~3.
\end{itemize}

% ============================================================
% References -- excluded from 2-page limit per rubric
% ============================================================
\clearpage

\refsheading
\begingroup
\small
\begin{enumerate}[label={[\arabic*]},leftmargin=*,itemsep=3pt,
                  topsep=2pt,parsep=0pt]
  \item A.~Vaswani, N.~Shazeer, N.~Parmar, J.~Uszkoreit, L.~Jones,
    A.\,N.~Gomez, L.~Kaiser, and I.~Polosukhin.
    \textit{Attention Is All You Need.}
    NeurIPS, 2017. arXiv:1706.03762.
  \item Y.-H.~Huang.
    \textit{attention-is-all-you-need-pytorch} (PyTorch codebase).
    GitHub, 2019. \texttt{github.com/jadore801120/attention-is-all-you-need-pytorch}
  \item D.~Elliott, S.~Frank, K.~Sima'an, and L.~Specia.
    Multi30K: Multilingual English-German Image Descriptions.
    \textit{ACL Workshop on Vision and Language (VL16)}, 2016.
  \item M.~Honnibal, I.~Montani, S.~Van~Landeghem, and A.~Boyd.
    \textit{spaCy: Industrial-strength NLP in Python.}
    Explosion AI, 2020. \texttt{spacy.io}
  \item PyTorch Development Team.
    \textit{PyTorch: High-Performance Deep Learning Library.}
    Version 2.x, 2023. \texttt{pytorch.org}
  \item R.~Sennrich, B.~Haddow, and A.~Birch.
    Neural Machine Translation of Rare Words with Subword Units.
    \textit{ACL}, 2016. (BPE utilities used in preprocessing.)
\end{enumerate}
\endgroup

\end{document}
